% Part: sets-functions-relations
% Chapter: relations
% Section: equivalence-relations

\documentclass[../../../include/open-logic-section]{subfiles}

\begin{document}

\olfileid{sfr}{rel}{eqv}

\olsection{ہم ارزیت دے تعلق}

اک سیٹ اُتے اکو ہون دا تعلق انعکاسی، متناظر تے متعدی ہوندا اے۔ ایہو
جہے تعلق $R$، جنہاں وچ ایہ تِنّے خاصیتاں ہون، بہت عام نیں۔

\begin{defn}[ہم ارزیت دا تعلق]
اک تعلق $R \subseteq A^2$ جیہڑا انعکاسی، متناظر تے متعدی ہووے،
\emph{ہم ارزیت دا تعلق} آکھیا جاندا اے۔ !!^{element}s $x$ تے $y$، جیہڑے
$A$ وچ نیں، \emph{$R$ دے لحاظ نال ہم ارز} آکھے جاندے نیں جے $Rxy$ ہووے۔
\end{defn}

ہم ارزیت دے تعلقاں توں \emph{ہم ارزیت دے طبقے} دا تصور نکلدا اے۔ ہم ارزیت
دا تعلق دائرے نوں وکھ وکھ حصیاں وچ ``ونڈ دیندا'' اے۔ ہر حصے دے اندر
ساریاں چیزاں اک دوجی نال تعلق رکھدیاں نیں؛ تے وکھرے حصیاں دیاں چیزاں
وچکار ایہ تعلق نہیں ہوندا۔ کدی کدی ایہناں حصیاں بارے \emph{سدھی} گل
کرنا فائدے مند ہوندا اے۔ ایس مقصد لئی اسی اک تعریف دیندے آں:

\begin{defn}\ollabel{def:equivalenceclass}
فرض کرو کہ $R \subseteq A^2$ اک ہم ارزیت دا تعلق اے۔ ہر $x \in A$ لئی،
$x$ دا $A$ وچ \emph{ہم ارزیت دا طبقہ} اوہ سیٹ $\equivrep{x}{R}
= \Setabs{y \in A}{Rxy}$ اے۔ سیٹ $A$ دا تعلق $R$ تحت \emph{خارج قسمت سیٹ}
$\equivclass{A}{R} = \Setabs{\equivrep{x}{R}}{x \in A}$ اے، یعنی ایہناں
ہم ارزیت دے طبقیاں دا سیٹ۔
\end{defn}

اگلا نتیجہ ہم ارزیت دے طبقے دی تعریف نوں درست ٹھہراؤندا اے: ایہ ثابت
کردا اے کہ ہم ارزیت دے طبقے واقعی $A$ دی ونڈ دے حصے نیں:

\begin{prop}
جے $R \subseteq A^2$ ہم ارزیت دا تعلق ہووے، تاں $Rxy$ اُدوں تے صرف اُدوں
جدوں $\equivrep{x}{R} = \equivrep{y}{R}$۔
\end{prop}

\begin{proof}
کھبے توں سجے پاسے دی سمت لئی، فرض کرو کہ $Rxy$، تے $z \in
\equivrep{x}{R}$ لو۔ تعریف مطابق فیر $Rxz$ اے۔ کیوں جو $R$ ہم ارزیت دا
تعلق اے، ایس لئی $Ryz$ اے۔ (کھول کے دسیے: کیوں جو $Rxy$ اے تے $R$
متناظر اے، ساڈے کول $Ryx$ اے؛ تے کیوں جو $Rxz$ اے تے $R$ متعدی اے،
ساڈے کول $Ryz$ اے۔) سو $z \in \equivrep{y}{R}$۔ ایس گل نوں عام کر کے،
$\equivrep{x}{R} \subseteq \equivrep{y}{R}$۔ پر بالکل ایسے طرح
$\equivrep{y}{R} \subseteq \equivrep{x}{R}$ وی اے۔ سو عنصراں نال تعیّن
مطابق $\equivrep{x}{R} = \equivrep{y}{R}$۔

سجے توں کھبے پاسے دی سمت لئی، فرض کرو کہ $\equivrep{x}{R} =
\equivrep{y}{R}$۔ کیوں جو $R$ انعکاسی اے، ایس لئی $Ryy$، سو $y \in
\equivrep{y}{R}$۔ ایس توں $y \in \equivrep{x}{R}$ وی اے، کیوں جو اسی
$\equivrep{x}{R} = \equivrep{y}{R}$ فرض کیتا اے۔ سو $Rxy$۔
\end{proof}

\begin{ex}
ہم ارزیت دے تعلقاں دی اک چنگی مثال باقی اُتے مبنی حساب توں ملدی اے۔
کسے وی $a$، $b$ تے $n \in \PosInt$ لئی، اسی $a \equiv_n b$ اُدوں تے
صرف اُدوں آکھدے آں جدوں $a$ نوں $n$ نال ونڈن توں اوہی باقی بچے جیہڑا
$b$ نوں $n$ نال ونڈن توں بچدا اے۔ (کجھ ہور علامتاں نال: $a \equiv_n b$
اُدوں تے صرف اُدوں جدوں کسے $k \in \Int$ لئی $a - b = kn$ ہووے۔) ہن
$\equiv_n$ ہر $n$ لئی ہم ارزیت دا تعلق اے۔ تے ٹھیک $n$ وکھرے ہم ارزیت دے
طبقے تعلق $\equiv_n$ توں بن دے نیں؛ یعنی $\equivclass{\Nat}{\equiv_n}$ دے
$n$ !!{element}s نیں۔ ایہ اوہ نیں: اوہناں عدداں دا سیٹ جیہڑے $n$ نال
بغیر باقی دے ونڈے جان، یعنی $\equivrep{0}{\equiv_n}$؛ اوہناں عدداں دا سیٹ
جیہناں نوں $n$ نال ونڈن اُتے باقی $1$ بچے، یعنی $\equivrep{1}{\equiv_n}$؛
\ldots؛ تے اوہناں عدداں دا سیٹ جیہناں نوں $n$ نال ونڈن اُتے باقی $n-1$
بچے، یعنی $\equivrep{n-1}{\equiv_n}$۔
\end{ex}

\begin{prob}
وکھاؤ کہ $\equiv_n$ ہر $n \in \PosInt$ لئی ہم ارزیت دا تعلق اے، تے
$\equivclass{\Nat}{\equiv_n}$ دے ٹھیک $n$ رکن نیں۔
\end{prob}

\end{document}
